\chapter{Reading a CT Scan}

\label{ch:read_ct}

Computed Tomography (CT) provides detailed cross-sectional imaging in multiple planes (axial, sagittal, coronal). While interpreting a CT scan can be daunting due to the volume of slices, applying a systematic approach and understanding the physical principles of tissue attenuation allows clinicians to identify major pathologies quickly and accurately.

---

\section*{1. CT Principles: Attenuation and Hounsfield Units (HU)}
CT images are constructed based on how much X-ray beam is absorbed (attenuated) by different tissues. The density of tissue is quantified using the \textbf{Hounsfield Unit (HU)} scale, calibrated against water:

\begin{center}
\begin{longtable}{@{}lll@{}}
\toprule
\textbf{Tissue/Substance} & \textbf{HU Range} & \textbf{Appearance on Standard Window} \\
\midrule
Acute Blood & $+55$ to $+80$ & Bright white (hyperdense) \\
Air & $-1000$ & Black (lowest density) \\
Bone / Calcification & $+700$ to $+3000$ & Bright white (highest density) \\
Cerebrospinal Fluid (CSF) & $0$ to $+15$ & Grey \\
Contrast Agent & $+100$ to $+300$ & Bright white \\
Fat & $-120$ to $-90$ & Dark grey \\
Soft Tissue / Muscle & $+30$ to $+60$ & Grey \\
Water & $0$ & Grey \\
\bottomrule
\end{longtable}
\end{center}

\subsection*{The Importance of Windowing}
The human eye can only distinguish about 30 shades of grey, whereas a CT scan captures thousands of HU. To optimize contrast, the computer groups HU ranges into "windows" defined by a \textbf{Window Width (WW)} (the range of HU shown) and a \textbf{Window Level (WL)} (the center of the HU range):
\begin{itemize}
    \item \textbf{Brain Window} (WW 80, WL 40): Narrow range to maximize contrast between grey and white matter and highlight acute hemorrhage.
    \item \textbf{Bone Window} (WW 2500, WL 600): Wide range to visualize cortical bone detail and fractures.
    \item \textbf{Lung Window} (WW 1500, WL -600): Extremely wide range to view air-filled lung parenchyma and bronchi.
    \item \textbf{Mediastinal / Soft Tissue Window} (WW 350, WL 40): Optimizes soft tissue visualization in the chest and abdomen.
\end{itemize}

---

\section*{2. Reading a Head CT: "Blood Can Be Very Bad"}
When interpreting a non-contrast head CT (NCCT), use this systematic mnemonic:

\subsection*{B --- Blood}
Look for acute hemorrhage, which appears bright white (hyperdense, +50 to +80 HU). Determine its location:
\begin{itemize}
    \item \textbf{Epidural Hemorrhage (EDH)}: Lenticular (biconvex/lemon-shaped) shape. Does not cross cranial sutures. Typically caused by middle meningeal artery rupture.
    \item \textbf{Subdural Hemorrhage (SDH)}: Crescent (banana-shaped) shape. Can cross sutures but does not cross dural attachments. Often caused by bridging vein rupture.
    \item \textbf{Subarachnoid Hemorrhage (SAH)}: Blood in the cisterns, sulci, and fissures. Often referred to as a "white star" in the suprasellar cistern. Usually caused by ruptured aneurysm.
    \item \textbf{Intraparenchymal Hemorrhage (IPH)}: Blood within the brain tissue itself (e.g., basal ganglia hemorrhage due to chronic hypertension).
\end{itemize}

\subsection*{C --- Canals (Cisterns)}
Inspect the fluid-filled subarachnoid spaces (cisterns) around the brainstem (suprasellar, quadrigeminal, prepontine). If these cisterns are compressed or obliterated (effacement), it indicates elevated intracranial pressure and potential brain herniation.

\subsection*{B --- Brain}
Compare the hemispheres side-to-side:
\begin{itemize}
    \item \textbf{Grey-White Differentiation}: Acute ischemic stroke causes cytotoxic edema, leading to loss of grey-white differentiation (e.g., "loss of the insular ribbon sign" or obscuration of the lentiform nucleus).
    \item \textbf{Midline Shift}: Draw a line from the anterior crista galli to the posterior internal occipital protuberance. Measure any lateral displacement of midline structures (like the septum pellucidum). Shift > 5 mm is a critical threshold.
    \item \textbf{Sulci}: Effacement (flattening) of sulci indicates diffuse cerebral edema or localized mass effect.
\end{itemize}

\subsection*{V --- Ventricles}
Examine the lateral, third, and fourth ventricles. Look for:
\begin{itemize}
    \item \textbf{Hydrocephalus}: Enlargement of the ventricles (dilation of temporal horns is an early sign).
    \item \textbf{Intraventricular Hemorrhage}: Blood pooling in the dependent posterior horns of the lateral ventricles.
\end{itemize}

\subsection*{B --- Bone (and Sinuses)}
Switch to a bone window. Check for skull fractures (linear, depressed, or basilar skull fractures). Check for fluid levels in the paranasal sinuses or mastoid air cells, which can indicate trauma (hemosinus) or infection.

---

\section*{3. Reading a Chest CT: Inside Out}
Examine the chest CT systematically in three planes using the soft tissue and lung windows:

\begin{itemize}
    \item \textbf{Mediastinum and Heart}: Evaluate the heart size, pericardium (for fluid), and mediastinal lymph nodes (abnormal if short-axis diameter > 10 mm). Trace the aorta (normal thoracic aorta < 4 cm; check for dissection flaps).
    \item \textbf{Pulmonary Vasculature}: On a contrast-enhanced CT Pulmonary Angiogram (CTPA), trace the main, lobar, and segmental pulmonary arteries. A dark filling defect surrounded by bright contrast indicates a \textbf{pulmonary embolism (PE)}.
    \item \textbf{Lung Parenchyma (Lung Window)}: Look for consolidation, ground-glass opacities (partial alveolar filling), nodules, cavities, or emphysematous destruction.
    \item \textbf{Pleural Space}: Look for pneumothorax (anteriorly in supine patients) and pleural effusions (posteriorly).
\end{itemize}

---

\section*{4. Reading an Abdomen/Pelvis CT: Organ by Organ}
Use a systematic trace through the abdomen:

\begin{itemize}
    \item \textbf{Solid Organs}:
    \begin{itemize}
        \item \textbf{Liver}: Check for homogeneous attenuation (fatty liver shows decreased attenuation, < 40 HU or less than spleen). Look for focal lesions (cysts, hemangiomas, metastases).
        \item \textbf{Gallbladder}: Look for wall thickening (> 3 mm), gallstones (calcified), or pericholecystic fluid.
        \item \textbf{Pancreas}: Check for swelling or fat stranding (acute pancreatitis) or mass lesions.
        \item \textbf{Spleen}: Normal size is < 12 cm. Look for lacerations or subcapsular hematomas in trauma.
        \item \textbf{Kidneys \& Adrenals}: Check for hydronephrosis, renal calculi (highly hyperdense on NCCT), or renal masses.
    \end{itemize}
    \item \textbf{Gastrointestinal Tract}: Trace from stomach to rectum. Look for bowel obstruction (dilated loops with air-fluid levels), bowel wall thickening, fat stranding (e.g., right lower quadrant stranding indicating appendicitis, or left lower quadrant stranding indicating diverticulitis).
    \item \textbf{Peritoneum and Space}: Search for free fluid (ascites or hemoperitoneum) in the hepato-renal recess (Morison's pouch) or pelvis. Search for free air (pneumoperitoneum) which pools anteriorly under the abdominal wall.
    \item \textbf{Aorta and Retroperitoneum}: Verify abdominal aortic diameter (< 3 cm is normal; > 3 cm represents an aneurysm).
\end{itemize}

---

\section*{5. Common Artifacts on CT}
\begin{itemize}
    \item \textbf{Beam Hardening}: Appears as dark bands or streaks between highly dense objects (e.g., bone or metal dental work). Caused by the low-energy photons of the X-ray beam being absorbed first, leaving only high-energy photons that distort the reconstruction.
    \item \textbf{Motion Artifact}: Appears as double borders or blurring. Caused by breathing or patient movement during acquisition.
\end{itemize}

---

\section*{6. Structured Reporting Example: Acute Appendicitis}

\begin{testbox}[Indications:]
24-year-old male with acute right lower quadrant pain, fever, and leukocytosis. Rule out appendicitis.
\end{testbox}
\begin{refbox}[Findings:]
\textbf{Solid Organs}: The liver, spleen, pancreas, gallbladder, and adrenal glands are normal in size and attenuation. No gallstones or biliary ductal dilatation. The kidneys enhance symmetrically; no hydronephrosis or perinephric fluid.\\
\textbf{Gastrointestinal Tract}: The stomach, duodenum, and small bowel loop sizes are within normal limits. In the right lower quadrant, the appendix is dilated, measuring 9 mm in outer diameter. There is circumferential thickening and enhancement of the appendiceal wall (target sign), accompanied by moderate surrounding fat stranding and a small fluid collection. A calcified appendicolith is noted at the base of the appendix. The colon is otherwise normal in appearance.\\
\textbf{Vessels and Peritoneum}: The abdominal aorta is normal. There is a small amount of free fluid localized in the right paracolic gutter and pelvis. No free air detected.\\
\textbf{Bones \& Pelvis}: No acute skeletal abnormalities.
\end{refbox}
\begin{pearlbox}[Impression:]
CT findings are indicative of acute appendicitis with an appendicolith and localized peritonitis. No findings of abscess perforation.
\end{pearlbox}

\clearpage
